% ---------------------------------------------------------------------
%  CHƯƠNG 4 — CÀI ĐẶT VÀ THỬ NGHIỆM
% ---------------------------------------------------------------------
\chapter{Cài đặt và thử nghiệm}
\label{ch:caidat}

\begin{muctieu}
Chương này mô tả môi trường triển khai, các thành phần chính kèm mã nguồn
minh hoạ, và giao diện demo của hệ thống gợi ý phòng trọ.
\end{muctieu}

\section{Môi trường triển khai}

Hệ thống được phát triển và thử nghiệm trên máy tính cá nhân với CPU Intel
Core i5, RAM 8~GB, hệ điều hành Windows 11 và Python 3.11 cùng các thư viện
đã nêu ở Bảng~\ref{tab:congcu}. Mã nguồn được quản lý bằng \code{git} và đặt
tại kho \url{https://github.com/tencuaem/tim-phong-tro-da-nang}.

\section{Các thành phần chính}

Mã nguồn~\ref{lst:featurize} minh hoạ bước chuẩn hoá và vec-tơ hoá đặc
trưng phòng trọ.

\begin{lstlisting}[language=Python, caption={Chuẩn hoá và vec-tơ hoá đặc trưng phòng trọ.}, label={lst:featurize}]
from sklearn.preprocessing import MinMaxScaler, OneHotEncoder
from sklearn.feature_extraction.text import TfidfVectorizer

# Normalize numeric features to the same [0, 1] scale
price_area_scaler = MinMaxScaler()
price_area = price_area_scaler.fit_transform(df[["gia_thue", "dien_tich"]])

# One-hot encode categorical district and amenity list
district_encoder = OneHotEncoder(sparse_output=False)
district_onehot = district_encoder.fit_transform(df[["quan"]])

# TF-IDF for the free-text description field
tfidf = TfidfVectorizer(max_features=500)
desc_vectors = tfidf.fit_transform(df["mo_ta"])
\end{lstlisting}

Mã nguồn~\ref{lst:score} minh hoạ hàm tính điểm gợi ý theo
Công thức~\eqref{eq:score} và xếp hạng phòng trọ.

\begin{lstlisting}[language=Python, caption={Tính điểm gợi ý và xếp hạng phòng trọ.}, label={lst:score}]
from sklearn.metrics.pairwise import cosine_similarity

# Weighted score across the four feature groups (see Equation 2.1)
def score_rooms(query_vec, room_vecs, weights):
    sim_price = cosine_similarity(query_vec["price_area"], room_vecs["price_area"])
    sim_area = cosine_similarity(query_vec["district"], room_vecs["district"])
    sim_amenity = cosine_similarity(query_vec["amenity"], room_vecs["amenity"])
    sim_desc = cosine_similarity(query_vec["desc"], room_vecs["desc"])

    w1, w2, w3, w4 = weights
    return w1 * sim_price + w2 * sim_area + w3 * sim_amenity + w4 * sim_desc

def top_k_rooms(scores, room_ids, k=5):
    # Return the k highest-scoring room ids, descending order
    ranked = sorted(zip(room_ids, scores), key=lambda x: x[1], reverse=True)
    return ranked[:k]
\end{lstlisting}

\section{Giao diện và demo}

\begin{figure}[H]
\centering
% Bỏ dấu % ở dòng dưới và thay ten-anh.png bằng ảnh thật của em
% \includegraphics[width=0.8\textwidth]{ten-anh.png}
\fbox{\begin{minipage}[c][5cm][c]{0.8\textwidth}\centering
  \itshape\color{gray500}[Chèn ảnh chụp màn hình giao diện tìm phòng trọ tại đây]
\end{minipage}}
\caption{Giao diện web cho phép nhập tiêu chí và xem danh sách phòng trọ
gợi ý kèm bản đồ vị trí.}
\label{fig:demo}
\end{figure}

\section{Kiểm thử}

Hệ thống được kiểm thử với ba nhóm ca. Nhóm một là tiêu chí thông thường
với ngân sách và khu vực rõ ràng. Nhóm hai là tiêu chí hẹp, ví dụ ngân sách
rất thấp kết hợp yêu cầu nhiều tiện ích, để kiểm tra hệ thống xử lý thế nào
khi không có phòng trọ nào thoả mãn đầy đủ. Nhóm ba là dữ liệu biên gồm
không nhập tiêu chí nào, nhập ngân sách âm và nhập khu vực không tồn tại
trong dữ liệu. Ở nhóm ba, hệ thống ban đầu báo lỗi khi ngân sách âm; lỗi
được khắc phục bằng cách kiểm tra và chặn giá trị không hợp lệ trước khi
đưa vào Công thức~\eqref{eq:score}.
